% ============================================================================
% Build: tectonic main.tex   (or: latexmk -pdf main.tex)
% Venue-neutral source. If a venue ever requires pdflatex + a shipped .bbl
% (arXiv does), both are produced by: tectonic --keep-intermediates main.tex
% Numbers in this paper are reproduced by commands in figures/README.md.
%
% FRAMING NOTE (2026-09-19): this paper deliberately makes NO novelty claim
% about the method. A prior-art scan (paper/prior-art/FINDINGS.md) established
% that verification-gated feature recovery is anticipated by Gupta & Nau (1995),
% ships in SolidWorks FeatureWorks (~1998) and PTC Creo (~2007), and is
% disclosed in expired patents from 2000. The contribution claimed here is
% availability, an honest evaluation, and a negative result. Do not reintroduce
% novelty language.
% ============================================================================
\documentclass[11pt]{article}

\usepackage[margin=1.1in]{geometry}
\usepackage[T1]{fontenc}
\usepackage[utf8]{inputenc}
\usepackage{lmodern}
\usepackage{microtype}
\usepackage{amsmath}
\usepackage{booktabs}
\usepackage{graphicx}
\usepackage{siunitx}
\usepackage{wasysym}
\usepackage{tikz}
\usetikzlibrary{positioning,arrows.meta}
\usepackage[numbers,sort&compress]{natbib}
\usepackage[hidelinks]{hyperref}
\usepackage{xspace}

\newcommand{\code}[1]{\texttt{\small #1}}
\newcommand{\eg}{e.g.\xspace}
\newcommand{\VERIFIED}{\textsc{verified}}
\newcommand{\PARTIAL}{\textsc{partial}}
\newcommand{\REFUSED}{\textsc{refused}}

\title{\textbf{\code{featuretree}: Verified Recovery of Editable CAD\\
Feature Trees from STEP, for Open CAD Systems}}
\author{Dan Newcome\\[2pt]
  \small punkfab\\
  \small\url{https://github.com/punkfab/featuretree}}
\date{September 2026}

\begin{document}
\maketitle

\begin{abstract}
Importing a STEP file into a parametric CAD system yields a single frozen solid:
the design history is gone, and no pad depth or hole diameter can be edited.
Recovering that history is a well-studied problem --- the acceptance criterion
we use was stated as a definition by \citet{TODO-gupta-nau} in 1995, and
commercial implementations have shipped since roughly 1998. It is nonetheless
absent from every open-source CAD system and from every callable library we
could find: the capability exists only inside proprietary seats costing
thousands of dollars per year, with no API. We describe \code{featuretree}, an
MIT-licensed, training-free recogniser that recovers an editable feature tree
from an exact B-rep, gates every recovery by re-executing the tree and comparing
the result against the input, and emits into FreeCAD and Onshape --- two systems
with no import-time feature recognition at all. On the eleven-part NIST MBE PMI
conformance corpus it returns \VERIFIED{} for 3 parts, \PARTIAL{} for 5, and
explicitly refuses 3, producing \emph{no silently incorrect trees}; all three
\VERIFIED{} recoveries exceed \SI{99.68}{\percent} IoU against the input. We
also report a negative result about our own acceptance criterion: a scalar
volume comparison lets over-cut and uncut material cancel, understating the true
geometric discrepancy by \num{1.56}$\times$ on one part, which a two-sided IoU
test eliminates. We claim no methodological novelty; we claim an available,
honest implementation and a measured account of where it fails.
\end{abstract}

\section{Introduction}

A solid model in a neutral exchange format carries geometry but no design
history. Opened in a parametric CAD tool it is one frozen body. Recovering an
editable history from such a body is the classical problem of automated feature
recognition (AFR), and it is genuinely old: \citet{TODO-gupta-nau} defined a
recovered feature set as correct precisely when subtracting those features from
stock reproduces the part, and \citet{TODO-vandenbrande} built a rule-based
system that rejected invalid candidates and iterated to a complete
decomposition. Both date from the mid-1990s.

The capability also ships. SolidWorks FeatureWorks has recognised features on
imported bodies since approximately 1998, and PTC's Creo Feature Recognition
Tool, available since Pro/ENGINEER Wildfire~4.0, states in its documentation
that it ``validates the model to ensure that the new features do not modify its
geometry''~\citep{TODO-creo-frt}. Geomagic Design~X recovers a parametric tree
and transfers it into five different CAD systems. The ISO~10303 standard went
further still: Parts~55, 108, 111 and 112 standardise construction history,
parameters and constraints, rolled into AP203 edition~2 and
AP242~\citep{TODO-pratt-kim}.

\paragraph{So why can you not do this?} Every implementation above is a GUI
feature inside a proprietary seat, priced from thousands to tens of thousands of
dollars per year, and none exposes an API a program can call. The two CAD
systems most accessible to an independent developer have nothing: Onshape's
documentation states that files imported from another system ``do not have a
feature tree in Onshape'', and FreeCAD has carried feature recognition as an
open request since 2011, with maintainers reporting no capacity to implement it.
The standardised route fared no better --- \citet{TODO-pratt-kim} reported in
2006 that ``there are at present no commercial STEP translators making use of''
the construction-history capabilities, and we are not aware of published work
since that revisits the question. Twenty years later the practical situation is
unchanged.

This paper is a response to that gap, not to a gap in the methods literature.

\paragraph{Contributions.}
\begin{enumerate}
\item \textbf{An open implementation (\S\ref{sec:system}).} A training-free,
  MIT-licensed recogniser from exact B-rep to an editable feature tree, emitted
  natively into FreeCAD PartDesign and Onshape, and to a \code{build123d} solid,
  from one neutral intermediate representation. It is a library, not a GUI
  feature.
\item \textbf{An evaluation with an explicit refusal class
  (\S\ref{sec:results}).} On the NIST MBE PMI conformance corpus we report
  \VERIFIED{}/\PARTIAL{}/\REFUSED{} per part, with reconstruction IoU for every
  accepted recovery. We are not aware of another AFR evaluation that reports a
  refusal rate alongside a success rate.
\item \textbf{A negative result about the acceptance criterion
  (\S\ref{sec:cancellation}).} Over-cut and uncut material enter a volume
  difference with opposite signs. We measure the resulting cancellation and show
  why a two-sided test is required.
\end{enumerate}

\paragraph{What this paper does not claim.} It does not claim novelty for
verification-gated recovery, for neutral feature representations, for
multi-target emission, or for round-tripping GUI edits. Each is anticipated;
\S\ref{sec:related} says by whom. The claim is availability, measurement, and
honesty about failure.

\section{Related work, and what it anticipates}
\label{sec:related}

We state the anticipating work plainly, because the contribution here does not
depend on it being absent.

\paragraph{The acceptance criterion is not new.} \citet{TODO-gupta-nau} require
that subtracting the recovered features from stock reproduce the part exactly,
and return a binary verdict --- a part with no feasible plan ``is considered
unmachinable''. Ours is the same gate with a measured tolerance in place of set
identity. In the learning literature, \citet{TODO-willis-fusion} accept a
reconstruction when IoU${}=1$, noting explicitly that a recovered sequence need
not match the ground-truth sequence: geometric equivalence rather than intent
recovery, published in 2021.

\paragraph{Neutral feature representations are not new.} The macro-parametric
approach~\citep{TODO-mun-han} defines tool-agnostic modelling commands replayed
into a target system, and Proficiency's commercial Universal Product
Representation~\citep{TODO-rappoport} held both parametric and B-rep versions of
each feature and regenerated native trees in six CAD systems from around 2000.

\paragraph{Multi-target emission and GUI round-trip are not new.} The BYU Neutral
Parametric Canonical Form line~\citep{TODO-shumway} supports concurrent editing
across NX, CATIA and Creo, capturing edits made in each GUI and propagating them
through a neutral database. It keys on database GUIDs where we key on stable
feature names --- an implementation difference, not an inventive one.

\paragraph{Query-based references are not new.} \citet{TODO-cascaval} formalise
references as ``queries that select elements from the constructed geometry'',
exactly the alternative to stored kernel identifiers that we adopt, with
correctness criteria we do not attempt.

\paragraph{What we did not find.} We searched for, and did not find: any
open-source B-rep-to-feature-tree recogniser; any AFR evaluation reporting a
graded verdict with a measured geometric residual per part; and any published
analysis of why the standardised AP242 construction-history route was not
adopted. Absence in one search is weak evidence, and we report these as unfound
rather than as nonexistent.

\section{System}
\label{sec:system}

The intermediate representation is a list of named, ordered features ---
\code{sketch}, \code{pad}, \code{pocket}, \code{revolve}, \code{prism\_cut},
\code{fillet} --- as plain JSON. Geometry is selected by \emph{query},
re-resolved against live geometry at each build rather than by persisted kernel
identifiers, which is what lets one IR drive structurally different backends and
lets edits be matched back by name. Authored IR uses human-meaningful names;
recovered IR uses generated ones (\code{pocket0}, \code{cut2}), stable enough
for the round trip but carrying no design intent.

Recovery proceeds cross-sectionally. Candidate axes are identified along which
the solid is prismatic, or about which it is a body of revolution; the outline
and through-holes are taken from cross-sections sampled at several heights;
blind features are deferred rather than drilled speculatively, since
over-cutting cannot be undone; and when the base extrude does not verify, the
residual is carved feature by feature, each remaining lump itself recognised as
a profile extruded along its own axis. Because several axes can appear
prismatic, candidates are not scored heuristically --- one is generated per axis
and the verifier selects the winner.

\paragraph{The gate.} Let $S$ be the input solid, $T$ a candidate tree and
$E(T)$ its re-execution. We accept when the volume difference and a
rotation-tolerant bounding-box comparison both fall within tolerance, and report
\PARTIAL{} with the residual otherwise. When no axis yields a recognisable
extrude or body of revolution, the recogniser \REFUSED{}s outright rather than
returning a best guess. Section~\ref{sec:cancellation} shows why this particular
invariant is weaker than it looks.

\section{Evaluation}
\label{sec:results}

\paragraph{Corpus.} The NIST MBE PMI Validation and Conformance Testing
artifacts, ``AP203 geometry only'' variants: eleven parts (CTC-01--05,
FTC-06--11), pure B-rep with no PMI annotations. These are human-designed
conformance artifacts intended to be difficult. NIST's own documentation notes
that they are not error-free reference files and that conformance checkers
report syntax errors in some of them, so a refusal may reflect the input rather
than the recogniser.

\begin{table}[htbp]
\centering
\caption{Recovery over the eleven-part NIST corpus. IoU is computed after
rotation-aware co-registration, because recovery works in the part's own frame
and the recovered solid is generally rotated relative to the input. Regenerate
with \code{paper/figures/run\_corpus.py} and \code{paper/figures/iou\_corpus.py}.}
\label{tab:corpus}
\begin{tabular}{llS[table-format=2.2]S[table-format=3.0]S[table-format=2.1]S[table-format=3.3]}
\toprule
Part & Verdict & {$|\Delta V|/V$ (\%)} & {Features} & {Time (s)} & {IoU (\%)} \\
\midrule
CTC-01 & \VERIFIED  &  0.10 &  51 & 15.7 & 99.841 \\
CTC-02 & \REFUSED   & {--}  & {--}& 66.4 & {--} \\
CTC-03 & \PARTIAL   & 41.46 &  16 &  6.1 & {--} \\
CTC-04 & \PARTIAL   &  2.43 & 270 & 37.2 & {--} \\
CTC-05 & \PARTIAL   & 51.80 & 106 & 11.6 & {--} \\
FTC-06 & \PARTIAL   & 31.45 &  41 &  8.7 & {--} \\
FTC-07 & \REFUSED   & {--}  & {--}&  8.4 & {--} \\
FTC-08 & \PARTIAL   & 52.55 & 218 & 48.3 & {--} \\
FTC-09 & \VERIFIED  &  0.19 &  80 &  5.1 & 99.684 \\
FTC-10 & \REFUSED   & {--}  & {--}&  8.5 & {--} \\
FTC-11 & \VERIFIED  &  0.00 &   2 &  1.6 & 100.000 \\
\bottomrule
\end{tabular}
\end{table}

\paragraph{Results.} Three parts \VERIFIED{} (\SI{27}{\percent}), five
\PARTIAL{} (\SI{45}{\percent}), three \REFUSED{} (\SI{27}{\percent}). All three
refusals were the same explicit rejection --- no recognisable extrude or body of
revolution --- rather than a crash or a guess.

The result we consider most important is that \textbf{no incorrect tree was
accepted}. Every \VERIFIED{} recovery exceeds \SI{99.68}{\percent} IoU against
the input. A \SI{27}{\percent} acceptance rate with an explicit refusal class is
a different and more useful contract than a higher rate with no soundness
signal, and it is the property that makes the output safe to edit.

\paragraph{Scale and comparability.} This is eleven parts. It establishes that
the verdict contract holds on non-trivial standard geometry; it does not
establish a rate over a part distribution, and no general recovery rate should
be read into it. FTC-11 in particular is a two-feature revolve recovered in
\SI{1.6}{\second} and should carry little weight. We note explicitly that these
numbers are \emph{not} comparable to the \SI{99}{\percent}-plus accuracies
reported on MFCAD++ and similar benchmarks: those measure per-face
classification against ground-truth labels on procedurally generated data,
whereas this measures whether a recovered program rebuilds a human-designed
conformance part.

\section{A negative result: volume comparison admits cancellation}
\label{sec:cancellation}

The scalar $|\Delta V|/V$ conceals its own structure. On CTC-01 the recovered
solid is \emph{smaller} than the original ($\Delta V =
\SI{-14850}{\cubic\milli\metre}$), so the residual cannot be attributed to uncut
chamfers, which would make it larger. Co-registering the solids and decomposing
by inclusion--exclusion gives \SI{19036}{\cubic\milli\metre} of \emph{over-cut}
against \SI{4186}{\cubic\milli\metre} of \emph{uncut} material: opposite signs,
partially cancelling, so the total geometric discrepancy
(\SI{0.159}{\percent}) is \num{1.56} times the \SI{0.101}{\percent} the verifier
reports.

Nothing in a volume-plus-extent test bounds that cancellation factor. A
\VERIFIED{} verdict under it should be read as ``reproduces the part's volume
and extent to within $\epsilon$'', not ``is the part to within $\epsilon$''. A
two-sided test --- IoU, or a symmetric-difference bound --- removes the failure
mode entirely, because over-cut and uncut material both increase the symmetric
difference. IoU is additionally the metric the CAD-program-inference literature
already reports, so adopting it buys comparability at the same time. We report
IoU in Table~\ref{tab:corpus} for this reason, and recommend it as the
acceptance criterion in place of the volume test the implementation currently
ships.

\paragraph{A practical obstacle.} The decomposition had to be derived from the
union by inclusion--exclusion, because direct OpenCASCADE difference and
intersection between these operands returned degenerate results --- an empty
intersection for solids whose bounding boxes coincide exactly. Any stronger
verifier built on Booleans will meet the same problem.

\section{Limitations}
\label{sec:limits}

\begin{itemize}
\item \textbf{Eleven parts.} A demonstration of the contract, not a study of a
  rate. A corpus-scale evaluation remains the obvious next step.
\item \textbf{\PARTIAL{} failures are unquantified.} Our IoU harness pre-filters
  candidate orientations by bounding-box match, and a wrong recovery has a wrong
  bounding box, so the five \PARTIAL{} parts have no IoU. We can say that they
  failed and by how much volume, but not how wrong they are geometrically.
\item \textbf{Equivalence, not intent.} Verification certifies that $E(T)$
  matches $S$, not that $T$ is the designer's history. Recovered feature names
  are generated and carry no design intent.
\item \textbf{Edge blends are outside the vocabulary.} Fillets and chamfers are
  not swept profiles and are left as sub-tolerance residual, so a chamfered part
  can verify while the recovered tree has sharp edges.
\item \textbf{Recovery is subtractive only.} Additive bosses cannot be
  recovered; such parts are reported \PARTIAL{}, not approximated.
\item \textbf{Freeform geometry is rejected, not approximated.} Splines,
  ellipses, lofts and sweeps admit no faithful sketch-and-pad tree.
\item \textbf{Tolerances are uncalibrated.} $\epsilon$ was chosen to separate
  observed successes from observed failures on the parts tested. It is not
  derived from a manufacturing tolerance model.
\item \textbf{Refusals are not attributed.} We did not separate refusals caused
  by geometry outside the vocabulary from refusals caused by defects in the
  input files, which NIST documents as present.
\end{itemize}

\section{Conclusion}

The methods here are not new; their availability is. A capability described in
the literature in 1995, shipped commercially since 1998 and standardised in
AP242 is still absent from open CAD, and that gap is a practical problem for
anyone building on FreeCAD, Onshape or a code-CAD toolchain. We have shown that
a training-free implementation reaches a useful contract on standard
conformance geometry --- accepting a quarter of the corpus with no incorrect
acceptances, and refusing rather than guessing on the rest --- and we have
measured a real weakness in the acceptance criterion we began with, together
with the two-sided replacement that removes it.

\section*{Availability}

\code{featuretree} is MIT-licensed and has been publicly available since
2026-06-26 at \url{https://github.com/punkfab/featuretree}; the recognition
component dates from 2026-07-10. Every number reported here is reproduced by the
commands in \code{paper/figures/README.md}, verified against the repository
state on 2026-09-19. The prior-art scan supporting \S\ref{sec:related} is in
\code{paper/prior-art/FINDINGS.md}.

This release is archived at
\href{https://doi.org/10.5281/zenodo.22848722}{\texttt{10.5281/zenodo.22848722}} (version~0.1.0);
the version-independent concept DOI, which always resolves to the latest version, is
\href{https://doi.org/10.5281/zenodo.22848721}{\texttt{10.5281/zenodo.22848721}}.

\bibliographystyle{plainnat}
\bibliography{references}

\end{document}
